% Options for packages loaded elsewhere
\PassOptionsToPackage{unicode}{hyperref}
\PassOptionsToPackage{hyphens}{url}
%
\documentclass[
]{article}
\usepackage{amsmath,amssymb}
\usepackage{iftex}
\ifPDFTeX
  \usepackage[T1]{fontenc}
  \usepackage[utf8]{inputenc}
  \usepackage{textcomp} % provide euro and other symbols
\else % if luatex or xetex
  \usepackage{unicode-math} % this also loads fontspec
  \defaultfontfeatures{Scale=MatchLowercase}
  \defaultfontfeatures[\rmfamily]{Ligatures=TeX,Scale=1}
\fi
\usepackage{lmodern}
\ifPDFTeX\else
  % xetex/luatex font selection
  \setmainfont[]{DejaVu Serif}
\fi
% Use upquote if available, for straight quotes in verbatim environments
\IfFileExists{upquote.sty}{\usepackage{upquote}}{}
\IfFileExists{microtype.sty}{% use microtype if available
  \usepackage[]{microtype}
  \UseMicrotypeSet[protrusion]{basicmath} % disable protrusion for tt fonts
}{}
\makeatletter
\@ifundefined{KOMAClassName}{% if non-KOMA class
  \IfFileExists{parskip.sty}{%
    \usepackage{parskip}
  }{% else
    \setlength{\parindent}{0pt}
    \setlength{\parskip}{6pt plus 2pt minus 1pt}}
}{% if KOMA class
  \KOMAoptions{parskip=half}}
\makeatother
\usepackage{xcolor}
\usepackage[margin=0.85in]{geometry}
\usepackage{longtable,booktabs,array}
\usepackage{calc} % for calculating minipage widths
% Correct order of tables after \paragraph or \subparagraph
\usepackage{etoolbox}
\makeatletter
\patchcmd\longtable{\par}{\if@noskipsec\mbox{}\fi\par}{}{}
\makeatother
% Allow footnotes in longtable head/foot
\IfFileExists{footnotehyper.sty}{\usepackage{footnotehyper}}{\usepackage{footnote}}
\makesavenoteenv{longtable}
\setlength{\emergencystretch}{6em} % prevent overfull lines
\providecommand{\tightlist}{%
  \setlength{\itemsep}{0pt}\setlength{\parskip}{0pt}}
\setcounter{secnumdepth}{-\maxdimen} % remove section numbering
\ifLuaTeX
  \usepackage{selnolig}  % disable illegal ligatures
\fi
\usepackage{bookmark}
\IfFileExists{xurl.sty}{\usepackage{xurl}}{} % add URL line breaks if available
\urlstyle{same}
\hypersetup{
  hidelinks,
  pdfcreator={LaTeX via pandoc}}

\author{}
\date{}

\hypersetup{pdftitle={Signed direction blocks and exact edge general position in rank-two independence cubes},pdfauthor={Ryutaro Yonezu (Independent Researcher)},pdfsubject={Research No.79 v1.0 Preprint; independent AI audit conditional PASS; not human peer reviewed}}
\begin{document}

\section{Signed direction blocks and exact edge general position in
rank-two independence
cubes}\label{signed-direction-blocks-and-exact-edge-general-position-in-rank-two-independence-cubes}

Research No.79, v1.0 Preprint --- 2026-11-21 JST

Ryutaro Yonezu (Independent Researcher)

Independent AI mathematical audit of Theorems 1--6: CONDITIONAL PASS; no
defect found. Application Appendix A: separate limited independent AI
audit of model membership and lower constructions; original upper bounds
depend on Nos.75--78. No human peer review or formal proof.

NOVELTY INCONCLUSIVE / REFERENCES PARTIAL / NOT HUMAN PEER REVIEWED /
All rights reserved.

\subsection{Abstract}\label{abstract}

We give a shared structural framework for edge general position in
isometric hypercube subgraphs, including the Fibonacci-type graphs
studied in Research Nos.75--78. Two-arm exclusions imply a star forest,
balanced direction relations and an antipodal support restriction across
maximal two-connected direction blocks. For independence cubes of rank
at most two, we identify domination exactly with one plus the
vertex-cover number of the graph of legal pairs. A partition certificate
then yields exact edge general position numbers for infinite families,
including arbitrary disjoint unions of triangle-free paths and cycles,
certain bipartite graphs, and linear/cyclic spacing cubes in explicit
ranges. We additionally prove a strict domination-bound gap and an exact
formula for double stars with at least two leaves on each side. A finite
catalogue covers all labelled triangle-free legal-pair graphs on at most
six labels. The unrestricted spacing-cube formula and a single
derivation of all previous Fibonacci-type formulas are not claimed.
Finite tests do not serve as premises of the analytic theorems.

\subsection{1. Shared structural
framework}\label{shared-structural-framework}

Let G be a finite isometric induced subgraph of Q\_n.~An edge general
position set (EGP set) S contains no three distinct edges on one
geodesic; gp\_e(G) is the maximum cardinality. A coordinate direction
class Theta\_i contains all edges flipping bit i; B2 is the sum of the
largest two class sizes (pad with zero if necessary). Every geodesic
flips its differing coordinates exactly once. For any subset J of
directions, the union of its complete classes is EGP if and only if the
Hamming diameter of the projection of V(G) onto J is at most two: if
larger, a geodesic between witnessing vertices contains three selected
edges; conversely at most two selected directions are flipped on any
geodesic. In particular B2 is always a lower bound. In a down-set of
Boolean words, shortest paths through intersection show isometry. On
active coordinates and for \textbar J\textbar\textgreater=3, the
projection diameter condition holds precisely when no legal pair of ones
lies inside J: such a pair and a third active singleton have projected
distance three; otherwise the projection consists of zero and
singletons.

\textbf{Lemma 1 (projection exclusion).} If two selected edges of
directions i,j share center c, any selected edge of another direction l
has fixed bits x\_i XOR x\_j=c\_i XOR c\_j. Otherwise its projection on
i,j matches exactly one leaf of the two-arm star. Orient the third edge
so that its initial endpoint w has l-bit equal to c\_l. Starting at the
opposite leaf, traverse the two arms through c to the matching leaf z,
then a shortest z-w path, then the third edge. The intervening path
flips neither i,j nor l. All the coordinates flipped by the
concatenation are distinct, so it is a geodesic, a contradiction.

\textbf{Lemma 2 (star forest).} If S uses at least three directions, the
spanning selected subgraph is a disjoint union of stars and isolates. A
simple three-edge selected path with distinct directions would be a
geodesic. The only possible repetition pattern is i,j,i; the two
overlapping two-arm stars then have opposite parity in (i,j), which
excludes every selected third-direction edge by Lemma 1. There is
therefore no simple three-edge path. As G is bipartite, every connected
selected component is a star or isolate.

\textbf{Theorem 1 (signed blocks).} Suppose S uses at least three
directions and has a star of degree at least three. Form R on its active
directions, putting an edge ij if they co-occur at a selected star
center, labelled sigma\_ij=c\_i XOR c\_j. This is well defined by Lemma
1 and the existence of a third active direction. Every cycle has label
XOR zero: any cycle omitting an active direction is tested by a selected
edge of that direction, and the XOR telescopes. A cycle with a chord
reduces to its two shorter cycles. A chordless spanning cycle of length
\textgreater=4 is impossible because R has a triangle from a selected
star. A spanning triangle has its three labels from that star, so its
XOR is zero. Consequently each connected component admits a bit
potential p with sigma\_ij=p\_i XOR p\_j, by propagation along paths.

Let D be an inclusion-maximal vertex set containing a maximum star's
clique in R and inducing a two-vertex-connected graph; write
d=\textbar D\textbar\textgreater=3. For an exterior selected edge of
direction l not in D, all constraints of R{[}D{]} hold, so its
projection is p or its complement q. For a selected edge in direction l
in D, the constraints of the connected R{[}D-\{l\}{]} force its other
bits to match p or q, so it is a radial edge incident to p or q. An
outside direction has at most one neighbor in D in R: otherwise
adjoining that direction would preserve connectivity after deletion of
any single vertex, contrary to maximality. These conclusions follow
solely from isometry and EGP, not from any recurrence or density
formula. They apply to all four earlier Fibonacci-type models, which are
down-sets.

\subsection{2. Antipodal support and
capacity}\label{antipodal-support-and-capacity}

For any D of size d\textgreater=3 and any p,q complementary on D, call
vertices projecting to p or q cores. Available radial edges flip a D bit
from a core. Their other endpoints are private leaves in this support: a
one-bit neighbor of one antipode is not a one-bit neighbor of the other,
and the unchanged outside coordinates identify the unique core. Exterior
edges join cores; the support includes only these two types. Write C for
the number of cores, a(v) for the radial degree at a core v, L=sum\_v
a(v), and Z for the number of zero-radial-degree cores. The full radial
family is EGP. In a projected geodesic a core pattern, once left by a D
flip, cannot be revisited. If both antipodes are visited, all D flips
are between those visits and only the first and last can be radial; if
just one is visited, at most two edges incident with it can be radial.
Thus there are at most two radial edges on any geodesic.

For an S supported here with at least three active directions, orient
each selected star from its center to its leaves. Every selected
exterior edge has a core leaf without any selected radial arm. Charge it
to an omitted radial edge at that leaf if a(v)\textgreater0, or to that
leaf if a(v)=0. All charges are distinct, proving
\textbar S\textbar\textless=L+Z. In particular, if Z=0 the best possible
supported EGP size is L. If the maximum selected degree is \textless=d,
the star-forest vertex count also gives
\textbar S\textbar\textless=floor(d(C+L)/(d+1)). These bounds include
the signed-block support of Theorem 1. The frozen v0.2--v0.4 supplements
supply further conditional matching, packing and weighted bounds; those
are relaxations, not assertions of universal equality. They are not
needed for the analytic exact theorems below.

\subsection{3. Rank-two domination and a partition
certificate}\label{rank-two-domination-and-a-partition-certificate}

Let K be any simple triangle-free graph with n\textgreater=1 labels and
m edges. Define G\_K as the cube induced by the cliques of K
(equivalently, independent sets of its complement). Its vertices are
empty, n singletons and m pairs. Its edges are spokes empty--singleton
and incidence edges pair--singleton. Thus N=1+n+m and
\textbar Theta\_i\textbar=1+deg\_K(i). Write tau(K) for minimum
vertex-cover size, alpha(K)=n-tau(K) for independence number, and
gamma(G\_K) for domination number.

\textbf{Theorem 2.} gamma(G\_K)=1+tau(K), and
gp\_e(G\_K)\textless=max(B2,m+alpha(K)).

Proof. Empty together with the singleton vertices of a minimum vertex
cover dominates G\_K. For a dominating set of any size, let T be the
chosen singleton labels, U=V(K)\textbackslash{}T, and P its chosen pair vertices viewed
as edges of K. Each edge of K{[}U{]} must lie in P. Choosing T and one
endpoint per such edge yields a cover of size at most
\textbar T\textbar+\textbar E(K{[}U{]})\textbar. If empty is chosen or P
contains an extra edge beyond E(K{[}U{]}), the dominating set has at
least 1+tau vertices. Otherwise empty is not chosen and P=E(K{[}U{]}); T
is nonempty because empty needs a neighbor. If two edges of K{[}U{]}
share an endpoint, cover those two by one vertex and every remaining
edge by one, giving
tau\textless=\textbar T\textbar+\textbar P\textbar-1. In the remaining
case K{[}U{]} is a matching. Every u in U is incident with a chosen pair
because its singleton must be dominated. Choose t in T. Retain T\{t\},
and select one endpoint per matching edge, taking the endpoint adjacent
to t if one exists. Triangle-freeness prevents both endpoints being
adjacent to t. This covers all edges of K, including those incident to
t; the cover has size \textbar T\textbar+\textbar P\textbar-1. This
completes the domination proof.

By Lemma 2, an EGP set using \textgreater=3 directions has star
components and isolates; one center per component is a dominating set in
G\_K. If there are c components then
\textbar S\textbar=N-c\textless=N-gamma=m+alpha. Sets using at most two
directions have size \textless=B2. QED.

For an arbitrary D subset of V(K), let T=V(K)\textbackslash{}D and select every edge
with direction in D incident to a vertex whose projection onto D is
zero. This one-core-pattern radial family is EGP even if
\textbar D\textbar\textless3: a projected geodesic can only leave or
enter that pattern once each. Its size is
\textbar D\textbar+cut\_K(D,T), since it includes the
\textbar D\textbar{} spokes and one incidence edge per crossing pair.
Write e\_K(X)=\textbar E(K{[}X{]})\textbar{} and define

c(K)=min\_\{T subset V(K)\} (\textbar T\textbar+e\_K(T)+e\_K(V(K)\textbackslash{}T)).

Then tau(K)\textless=c(K) (T plus one endpoint per uncovered edge is a
cover), and the largest one-pattern radial family has n+m-c(K) edges.
Theorem 2 yields the exact sandwich

max(B2,n+m-c(K)) \textless= gp\_e(G\_K) \textless= max(B2,n+m-tau(K)).
(1)

\textbf{Theorem 3 (partition certificate).} If there is a partition into
T,D with T independent and \textbar T\textbar+e\_K(D)=tau(K), then

gp\_e(G\_K)=max(B2,m+alpha(K)). (2)

In fact c(K)=tau(K) if and only if such a certificate exists. For the
reverse direction, a partition cost equal to tau satisfies
tau\textless=\textbar T\textbar+e\_K(D)\textless=cost=tau, so e\_K(T)=0.
Necessarily K{[}D{]} is a matching plus isolates: two adjacent uncovered
edges permit the cover to save one vertex, contradicting minimality.
This is an exact criterion for the partition bound to reach the
domination term, not a necessary condition for the latter term to be
attained by arbitrary EGP sets.

\textbf{Corollary 3a.} If K is triangle-free and max degree \textless=2,
equation (2) holds. Its components are paths (including isolates) or
cycles. In a path or an even cycle choose the smaller alternating side
as T; it is independent, D is independent, and a matching of size
\textbar T\textbar{} proves minimum cover size \textbar T\textbar. In an
odd cycle of length 2r+1\textgreater=5, choose r alternate vertices as
T; D has exactly one edge. The cycle has cover number r+1 because an
independent set has at most r vertices (each selected vertex has a
distinct unselected successor) and the construction reaches r. Thus
\textbar T\textbar+e(D)=r+1. Combining these certificates over
components proves the result by additivity of vertex-cover size. More
generally, a disjoint union of bipartite components each having a
matching saturating its smaller part has a certificate given by the
union of those smaller parts. Both statements allow arbitrarily many
coordinate labels.

\subsection{4. Strict domination gap: an exact double-star
family}\label{strict-domination-gap-an-exact-double-star-family}

\textbf{Theorem 4.} Let K be a double star with adjacent centers a,b,
with p\textgreater=2 leaves adjacent only to a and q\textgreater=2
leaves adjacent only to b. Write s=p+q. Then gp\_e(G\_K)=2s. The
domination term is 2s+1, strictly too large.

Proof. K has n=s+2, m=s+1, tau=2 (the two centers cover K and disjoint
edges force at least two cover vertices), and B2=s+4\textless=2s.
Equation (1) gives gp\_e\textless=2s+1. Choose all leaves as D in the
radial construction, giving s spokes and s crossing incidence edges,
hence 2s edges.

It remains to exclude an EGP set of size 2s+1. Such a set needs
\textgreater=3 directions because B2\textless=2s. It is a star forest
and has exactly N-(2s+1)=3 components. Any choice of star centers
(isolates included) is a dominating triple. We show that the only
dominating triple is \{empty,\{a\},\{b\}\}. A dominating triple not
containing empty must contain at least one singleton to dominate empty.
If it has t\textgreater=1 singleton vertices, its remaining 3-t pair
vertices can dominate at most 2(3-t) singleton vertices, so at most
t+2(3-t)\textless=5 singleton vertices are dominated. There are
s+2\textgreater=6. If empty belongs to the triple, pair vertices must be
dominated by the other two selected vertices. If \{a\} is not chosen,
the p\textgreater=2 pairs consisting of a and an a-leaf each need their
own chosen leaf singleton or that pair vertex. Thus both remaining
positions must be used for these pairs (or p\textgreater2 is already
impossible). Neither of these positions can dominate any pair consisting
of b and a b-leaf; contradiction. By symmetry \{b\} must be chosen too.
The triple is unique.

Every noncenter vertex now has to attach to an adjacent center in its
star. Each leaf singleton must attach to empty, and each center--leaf
pair must attach to its corresponding center singleton. The bridge pair
\{a,b\} attaches either to \{a\} or \{b\}. In the former case, at \{a\}
the bridge direction b and an a-leaf direction i are selected arms, with
center parity zero on b,i. At \{b\} any b-leaf direction j is a third
selected direction; its fixed bits have b=1,i=0. This contradicts Lemma
1. Attachment to \{b\} is symmetric. No EGP set reaches the domination
term, and the constructed lower bound is exact. QED.

\subsection{5. Spacing cubes: exact ranges and
limits}\label{spacing-cubes-exact-ranges-and-limits}

Let G(n,h) be the cube of binary words of length n whose ones have
pairwise index distance \textgreater=h+1. Let A\_m\^{}(h)=1 for
m\textless=0 and A\_m=A\_(m-1)+A\_(m-h-1) for m\textgreater=1. Let
G\_c(n,h) use the analogous cyclic distance. These models are known
generalized Fibonacci/Lucas cubes; the run-avoidance model of earlier
Nos.77--78 differs.

\textbf{Theorem 5 (linear rank-two range).} For h\textgreater=2 and
2\textless=n\textless=2h+2, gp\_e(G(n,h))=A\_(n-1)\^{}(h).

If n\textless=h+1 the graph is an n-edge star. Otherwise put t=n-h-1,
with 1\textless=t\textless=h+1. Legal pairs connect left positions 1..t
and right positions h+2..n. These parts are disjoint, the graph is
bipartite (other labels isolated), and pairs (i,i+h+1) form a matching
of size t saturating both sides. Thus tau=t, m=t(t+1)/2,
N-gamma=n+m-t=A\_(n-1). Each coordinate class has at most 1+t edges.
B2\textless=2(1+t)\textless=n+m-t: the difference is h-1-2t+t(t+1)/2 =
h-2+(t-1)(t-2)/2, nonnegative for h\textgreater=2,t\textgreater=1.
Taking the first h+1 coordinate classes selects every word meeting those
positions, each exactly once, so yields A\_n-A\_(n-h-1)=A\_(n-1). This
also follows from the partition certificate. QED.

\textbf{Theorem 6 (cyclic rank-two range).} For h\textgreater=1,
gp\_e(G\_c(n,h))=n for 2\textless=n\textless=2h+2 and
gp\_e(G\_c(2h+3,h))=3h+4.

When n\textless=2h+1 the legal-pair graph is empty; at n=2h+2 it is a
perfect matching with tau=h+1. At n=2h+3 it is a cycle on n labels with
step h+1, since gcd(h+1,2h+3)=1. No three ones fit cyclically because
the sum of three gaps would have to be
\textgreater=3(h+1)\textgreater2h+3 for h\textgreater=1. Apply Corollary
3a: for the matching, m+alpha=n and B2=4\textless=n; for the cycle,
m+alpha=(2h+3)+(h+1)=3h+4 and B2=6\textless=3h+4. This includes the
previously known Lambda\_5 value 7, not a new claim for that instance.
QED.

The frozen v0.2 proof gives the analytic lower bound A\_(n-1) for every
h\textgreater=2,n\textgreater=2 and bounds all complete coordinate
unions and pure antipodal radial families by that number. The v0.3--v0.4
supplements report coincident computational bounds for 42 spacing
models: h=2,n=2..16; h=3,n=2..15; h=4,n=2..14. These finite results are
inherited working results, not necessary for Theorems 1--6; their
earlier verifiers have not been fully replayed in v0.6. The frozen v0.5
also identifies a barrier: for n\textgreater=2h+3 the domination number
is \textless=A\_(n-h-1)-1, so the domination bound alone cannot settle
the all-n linear conjecture. This is a limitation of the method, not a
counterexample to that conjecture.

\subsection{6. Finite catalogue and boundary of
scope}\label{finite-catalogue-and-boundary-of-scope}

The new verifier enumerates all labelled simple triangle-free graphs K
for 1\textless=n\textless=6 (6,228 in total, including isolated labels),
checks vertex-cover and partition costs by direct enumeration, and finds
6,138 with c=tau. Every remaining graph (90) is isomorphic to the double
star p=q=2: its degree multiset is (1,1,1,1,3,3), the degree-three
vertices adjacent, and it has five edges. These facts force the asserted
isomorphism. Theorem 4 covers all 90. Hence the catalogue is exhausted
by a finite classification and the analytic theorems. As supplementary
checks, the code computes actual BFS distances and forbidden triples by
all orientations/orders in 51 models on at most four labels and in the
remaining representative double star (optimum 8, 92 forbidden triples);
it also checks unique dominating triples for 25 double stars p,q=2..6.
This is author-side corroboration, not independent review. The infinite
theorems do not depend on this catalogue.

We deliberately stop generalization here. The paper provides a common
structural mechanism rather than a single fresh exact formula subsuming
Nos.75--78. The full linear spacing conjecture, the cyclic range
n\textgreater2h+3, arbitrary rank-two legal graphs not covered by the
certificate/double-star theorem, and higher-rank independence cubes
remain open within this work. No global novelty or priority claim is
made. Theorems 1--6 received a conditional independent AI mathematical
audit PASS. The unrestricted cases remain open.

\subsection{References (partial)}\label{references-partial}

\begin{itemize}
\item
  P. Manuel, R. Prabha, S. Klavžar, The edge general position problem,
  Bull. Malays. Math. Sci. Soc. 45 (2022), 2997--3009, DOI
  10.1007/s40840-022-01319-8. The two-Theta-class construction in
  partial cubes is prior work (Lemma 3.1).
\item
  J. Tian, S. Klavžar, E. Tan, Extremal Edge General Position Sets in
  Some Graphs, Graphs Combin. 40 (2024), article 40, DOI
  10.1007/s00373-024-02770-z.
\item
  P. Manuel, R. Prabha, S. Klavžar, Generalization of edge general
  position problem, Art Discrete Appl. Math. 8 (2025), \#P1.02, DOI
  10.26493/2590-9770.1745.5f4.
\item
  Z. Hamed-Labbafian, M. A. Henning, M. Tavakoli, Edge general position
  in graphs: Graph products, integer linear programming and some
  applications, Discrete Appl. Math. 386 (2026), 1--8, DOI
  10.1016/j.dam.2026.01.008. An ILP approach is prior work.
\item
  V. Ullas Chandran S. V., S. Klavžar, J. Tuite, The General Position
  Problem: A Survey, arXiv:2501.19385 (2025), Section 3.6.
\item
  S. Klavzar and E. Tan, Edge General Position Sets in Fibonacci and
  Lucas Cubes (2023), DOI 10.1007/s40840-023-01517-y,
  https://arxiv.org/abs/2304.10114. Prior two-class constructions and
  Fibonacci/Lucas EGP context.
\item
  P. Codara and O. M. D'Antona, Generalized Fibonacci and Lucas cubes
  arising from powers of paths and cycles, Discrete Mathematics 339
  (2016), 270--282, DOI 10.1016/j.disc.2015.08.012. Prior graph model
  and enumeration.
\item
  Author's earlier research Nos.75--78: related constructions and
  proofs. Their claimed exact formulas are not re-proved or claimed as
  new results here.
\end{itemize}

Targeted preliminary searches for edge general position with daisy
cubes, independence cubes, domination, and generalized Fibonacci cubes
did not identify a direct matching theorem. Search-result/abstract scope
only; full-text/citation survey remains incomplete. NOVELTY INCONCLUSIVE
/ REFERENCES PARTIAL.

\newpage

\section{Research No.79 v0.6 --- Application Appendix A:
Nos.75--78}\label{research-no.79-v0.6-application-appendix-a-nos.7578}

\textbf{Editorial Appendix A --- 2026-11-21 JST.} This appendix received
a separate limited independent AI audit of its down-set assumptions,
construction counts and selected finite cases. The earlier independent
audit of Theorems 1--6 did not cover it. Upper bounds and equality
claims remain dependent on Nos.75--78. NOVELTY INCONCLUSIVE / REFERENCES
PARTIAL / NOT HUMAN PEER REVIEWED / All rights reserved.

\subsection{A1. The exact distinction between applicability and a
unified
proof}\label{a1.-the-exact-distinction-between-applicability-and-a-unified-proof}

No.79, Lemmas 1--2 and Theorem 1 are shared \textbf{structural} results
for isometric induced subgraphs of a hypercube; the complete exact
formulas of Nos.75--78 do \textbf{not} follow from those structural
results alone. No.79, Theorems 2--6, are \textbf{rank-two} theorems and
cannot be applied automatically to every dimension of Nos.75--78.
Individual \emph{upper bounds and equality classifications} below are
imported from the respective earlier mathematical manuscripts, not
re-proved independently here. No claim of a new unified derivation of
all four formulas is made.

\subsection{A2. One proof of isometry for each of the four graph
families}\label{a2.-one-proof-of-isometry-for-each-of-the-four-graph-families}

Each family has as vertices a set of binary words closed under changing
any 1 to 0:

\begin{enumerate}
\def\labelenumi{\arabic{enumi}.}
\tightlist
\item
  No.75: length-\texttt{n} words avoiding adjacent ones,
  n\textgreater=2;
\item
  No.76: length-\texttt{n} cyclic words avoiding cyclically adjacent
  ones, n\textgreater=3;
\item
  No.77: length-\texttt{n} words avoiding \texttt{k} consecutive ones,
  k\textgreater=3 and n\textgreater=2;
\item
  No.78: length-\texttt{n} cyclic words avoiding cyclic runs of
  \texttt{k} ones, n\textgreater=k\textgreater=3.
\end{enumerate}

Deleting ones cannot create a forbidden all-ones pattern. For two
allowed words x,y, first delete each bit in x\textbackslash{}y, reaching x AND y; then
insert each bit in y\textbackslash{}x, remaining within the support of y. Every
intermediate word is legal. This path has Hamming length \textbar x XOR
y\textbar, so every family is an isometric induced subgraph of
Q\_n.~Thus for every EGP set S in any of the four models:

\begin{itemize}
\tightlist
\item
  No.79's two-arm projection exclusion, Lemma 1, applies.
\item
  If S uses \textgreater=3 coordinate directions, the selected subgraph
  is a star forest (Lemma 2).
\item
  If S uses \textgreater=3 directions and has a selected star of degree
  \textgreater=3, the signed direction relation, balanced potentials,
  maximal 2-connected block, and antipodal radial/exterior support
  conclusion of Theorem 1 apply \textbf{subject to its stated
  hypotheses}.
\item
  A complete union of at most two coordinate classes is always EGP.
\end{itemize}

This is a direct analytic check of every model's hypotheses, not just an
assertion based on titles.

\subsection{A3. No.75 --- linear Fibonacci cube
Gamma\_n}\label{a3.-no.75-linear-fibonacci-cube-gamma_n}

\textbf{Construction by the general two-direction principle.} For
n\textgreater=2, both endpoint classes have
\textbar Theta\_1\textbar=\textbar Theta\_n\textbar=F\_n, hence the
union is EGP with exactly 2F\_n edges. The size identity follows from
the legal word count on the remaining n-1 positions after imposing a 1
at the endpoint.

\textbf{Specific optimality still needed.} No.75 v0.10, §§2--5, supplies
the upper bound by (i) star-forest density for maximum degree
\textless=3, (ii) projection exclusion and Fibonacci base/recurrence
certificate for degree 4, and (iii) a five-arm exclusion count for
larger degrees. Its equality classification singles out the two boundary
coordinate classes. Theorem 1 of No.79 explains the same star/projection
mechanism, but does not by itself verify those numerical inequalities or
uniqueness.

\textbf{Bridge:} structural hypothesis PASS; lower construction PASS;
final upper bound/uniqueness CONDITIONAL ON No.75 proof, not newly
proved here.

\subsection{A4. No.76 --- cyclic Lucas cube Lambda\_n: the key
three-direction
application}\label{a4.-no.76-cyclic-lucas-cube-lambda_n-the-key-three-direction-application}

The coordinate classes have \textbar Theta\_i\textbar=F\_(n-1). For
n\textgreater=5, their largest two-class union has only 2F\_(n-1) edges,
strictly fewer than the No.76 optimum F\_n+F\_(n-2) by
F\_(n-4)\textgreater0. Therefore a proof that only considers complete
coordinate-class unions \textbf{cannot} explain this model.

\textbf{Construction by No.79 antipodal radial support.} Fix three
consecutive cyclic directions D=\{0,1,2\}. Select every direction-D edge
whose other two D bits are zero. Equivalently these are radial edges at
the 000 core in each fixed outside-coordinate fiber. No.79 §2 proves
that the full radial family is EGP. Counting the middle-direction edges
gives F\_(n-1); each outer direction contributes F\_(n-2), for a total

\begin{verbatim}
F_(n-1) + 2 F_(n-2) = F_n + F_(n-2).
\end{verbatim}

For n=3 this is the three-edge star; for n=4 it equals the two-class
bound; for n\textgreater=5 it strictly beats that bound. The selected
three-arm center with all other bits zero generates a signed relation
triangle on D and has potential 000, exactly as in No.79 Theorem 1.

\textbf{Specific optimality still needed.} No.76 v0.1--v0.3 supplies the
substantial cyclic gap and local star-pair capacity arguments, including
finite recurrence bases and the exceptional n=6 relaxation case. No.79's
structural theorem alone does not replace those checks. No.76 does not
claim full equality-set uniqueness.

\textbf{Bridge:} structural hypothesis PASS; explicit application of
No.79 radial family PASS; construction cardinality PASS; upper bound
CONDITIONAL ON No.76 proofs, not independently reproduced by this
appendix.

\subsection{A5. No.77 --- higher-order linear run-avoiding cube
Q\_n(1\^{}k)}\label{a5.-no.77-higher-order-linear-run-avoiding-cube-q_n1k}

For k\textgreater=3,n\textgreater=2, let A\_n\^{}(k) count the legal
words. The two endpoint classes have size A\_n-A\_(n-1) each. The No.79
two-class principle yields an EGP set of size

\begin{verbatim}
2 (A_n - A_(n-1)).
\end{verbatim}

All models are down-sets, so the generic star/projection/signed-block
framework applies to other candidate extremizers. The No.77 proof
additionally needs its density inequality B\textgreater=10A\_n/11 and
the six-separated-direction exclusion for large stars; these yield the
\textbf{strict} upper bound for any EGP set using \textgreater=3
directions. The original equality classification then gives the endpoint
union for n\textgreater=k+1 and any two complete classes for
2\textless=n\textless=k.

\textbf{Bridge:} structural hypothesis PASS; construction cardinality
PASS; final upper bound/equality CONDITIONAL ON No.77 proof.

\subsection{A6. No.78 --- higher-order cyclic run-avoiding cube
H(n,k)}\label{a6.-no.78-higher-order-cyclic-run-avoiding-cube-hnk}

For n\textgreater=k\textgreater=3, all coordinate classes have size
C\_n\textsuperscript{(k)-A\_(n-1)}(k). The No.79 two-class principle
yields

\begin{verbatim}
2 (C_n - A_(n-1))
\end{verbatim}

selected EGP edges for any pair of distinct directions. The
down-set/isometry, projection exclusion and star forest requirements
hold verbatim. The No.78 proof's missing case k=3 relies on a separate
transfer-matrix certificate (1,296 four-gap/center bases) and a finite
n=4 equality-classification computation. For k\textgreater=4 it uses the
original analytic density and excluded-projection arguments. The
uniqueness of maxima is imported from No.78, not deduced from No.79's
general structure alone.

\textbf{Bridge:} structural hypothesis PASS; construction cardinality
PASS; final upper bound/equality CONDITIONAL ON No.78 proof and its
numerical certificate.

\subsection{A7. Summary and closure
recommendation}\label{a7.-summary-and-closure-recommendation}

\begin{longtable}[]{@{}
  >{\raggedright\arraybackslash}p{(\columnwidth - 6\tabcolsep) * \real{0.2500}}
  >{\raggedright\arraybackslash}p{(\columnwidth - 6\tabcolsep) * \real{0.2500}}
  >{\raggedright\arraybackslash}p{(\columnwidth - 6\tabcolsep) * \real{0.2500}}
  >{\raggedright\arraybackslash}p{(\columnwidth - 6\tabcolsep) * \real{0.2500}}@{}}
\toprule\noalign{}
\begin{minipage}[b]{\linewidth}\raggedright
Model
\end{minipage} & \begin{minipage}[b]{\linewidth}\raggedright
Fits No.79 structural framework
\end{minipage} & \begin{minipage}[b]{\linewidth}\raggedright
A No.79 construction reaches its reported optimum
\end{minipage} & \begin{minipage}[b]{\linewidth}\raggedright
No.79 alone reproves upper bound
\end{minipage} \\
\midrule\noalign{}
\endhead
\bottomrule\noalign{}
\endlastfoot
No.75 Gamma\_n & Yes, all n\textgreater=2 & Yes, two complete endpoint
classes & \textbf{No} \\
No.76 Lambda\_n & Yes, all n\textgreater=3 & \textbf{Yes,
three-direction 000-core radial family} & \textbf{No} \\
No.77 Q\_n(1\^{}k) & Yes, k\textgreater=3,n\textgreater=2 & Yes, two
endpoint classes & \textbf{No} \\
No.78 H(n,k) & Yes, n\textgreater=k\textgreater=3 & Yes, any two
complete classes & \textbf{No} \\
\end{longtable}

\textbf{Proposed claim for the final paper:} ``Theorem 1 applies
individually to all four earlier run-avoiding and Fibonacci/Lucas cube
families; its two-direction and antipodal radial constructions recover
their corresponding lower-bound mechanisms. Their exact upper bounds and
extremizer classifications remain separate model-specific consequences
proved in Nos.75--78, not new corollaries of Theorem 1 alone.''

This is a stronger, accurately scoped application statement than merely
listing the four graph names. It does not justify advertising a complete
unified independent proof of all four prior theorems.

\subsection{A8. Local spot-check
scope}\label{a8.-local-spot-check-scope}

Run \texttt{python\ -B\ No79\_V06\_APPLICATION\_CHECK\_75\_78.py}. It
uses \textbf{only Python's standard library} and exhaustively enumerates
legal vertices and construction edges over 49 selected parameter
instances: 7 (No.75), 6 (No.76), 21 (No.77), and 15 (No.78). For each,
the verifier checks downward closure, exact construction count, and
computes the maximum number of selected edges on any shortest path by an
independently written shortest-path DAG DP. It does \textbf{not} check
arbitrary EGP upper bounds, all dimensional parameters, or uniqueness.
The analytic arguments above, not these finite tests, establish the
model membership and lower constructions for all stated parameters.

\subsection{A9. Source identity and editorial
discipline}\label{a9.-source-identity-and-editorial-discipline}

The consulted mathematical source texts, exact SHA256:

\begin{itemize}
\tightlist
\item
  \texttt{No79\_UNIFIED\_V06\_CLOSURE/MANUSCRIPT\_V06.md} ---
  \nolinkurl{eadbe8ca533cd7dfafac9dfad37781088af3f8e95f7140b785c8937e087123cf}
\item
  \texttt{No75\_V010/RESULTS.md} ---
  \nolinkurl{d9af8c385e3c67cec908a8dde821878b8d95805f207826d89282f572481a14aa}
\item
  \texttt{No76\_LUCAS\_V03/RESULTS.md} ---
  \nolinkurl{3aec6fd5cb5f77469743a136ade8c86d0f9decdac916f121bea4753b7bd9b092}
  (references nested v0.1/v0.2 originals)
\item
  \texttt{No77\_GENERAL\_K\_V01/MANUSCRIPT.md} ---
  \nolinkurl{83caf0c565d5f81e9adf96b4cff21245dc962353c8b4bb04f74fd7c68dea49a9}
\item
  \texttt{No78\_CYCLIC\_K\_V02/MANUSCRIPT\_V02.md} ---
  \nolinkurl{cee2fc12ba41b2569eaa39e5b6b1e1912f1d778e055883222faed6ed34708a9c}
  (references nested v0.1 original)
\end{itemize}

The above texts are the available frozen research-source materials
consulted for this appendix, not an independent cross-audit of the
currently public repository PDFs. No original byte or mathematics is
changed. A targeted limited application audit was performed for this
v1.0 preprint. It is not a full independent re-audit of Nos.75--78
or publication authorization.

\end{document}
